\documentclass{article}
\usepackage[a4paper,margin=2.5cm]{geometry}
\usepackage{amsmath}

\begin{document}

\section{Convolution}

\subsection{Causal convolution using quantics MPS}

We consider the causal convolution
\begin{equation}
    F(t)=\int_{T_0}^{t} K(t-v)S(v)\,dv,
    \qquad T_0\leq t<T_M,
\end{equation}
where $K$ is the convolution kernel and $S$ is the signal.

\paragraph{Discretization.}

Let $L=2^N$ be the number of grid points, and define
\begin{equation}
    \Delta t=\frac{T_M-T_0}{L}.
\end{equation}

The signal and the output are discretized on the same grid:
\begin{equation}
    t_m=T_0+m\Delta t,
    \qquad m=0,\ldots,L-1.
\end{equation}

The kernel is defined on the corresponding time-lag grid:
\begin{equation}
    \tau_k=k\Delta t,
    \qquad k=0,\ldots,L-1.
\end{equation}

Thus, the discrete functions are
\begin{equation}
    S_m=S(t_m),\qquad
    K_k=K(\tau_k),\qquad
    F_m=F(t_m).
\end{equation}

Scaling the inclusive convolution gives the following uncorrected quadrature
expression; its endpoint convention is discussed below:
\begin{equation}
    \boxed{
    F_m \approx \Delta t
    \sum_{j=0}^{m}K_{m-j}S_j.
    }
    \label{eq:causal_convolution}
\end{equation}

\paragraph{Endpoint convention and quadrature corrections.}

The current \texttt{conv} operation returns the inclusive discrete sum
\begin{equation}
    C_m=\sum_{j=0}^{m}K_{m-j}S_j,
\end{equation}
without the factor $\Delta t$. This is a discrete convolution, not yet a
particular continuous-integration rule. The interval $[T_0,t_m]$ contains
$m$ grid intervals but $m+1$ samples. For example, if $K=S=1$, the integral
is $m\Delta t$, whereas the uncorrected inclusive sum gives $(m+1)\Delta t$.
In particular, the uncorrected value at $m=0$ need not vanish.

To reproduce the previous exclusive-prefix convention in the time-lag
coordinate, use the right-rectangle rule in the physical signal time $v$:
\begin{equation}
    \boxed{
    F_m^{\mathrm{R}}=
    \Delta t\sum_{j=1}^{m}K_{m-j}S_j
    =\Delta t\left(C_m-S_0K_m\right).
    }
    \label{eq:convolution_right_rectangle}
\end{equation}
Indeed, writing $k=m-j$ gives an exclusive sum over lag indices
$k=0,\ldots,m-1$. Consequently:
\begin{itemize}
    \item For an adiabatic activation $S(v)=\lambda(v)$ with
    $S_0=\lambda(T_0)=0$, the correction vanishes:
    $F_m^{\mathrm{R}}=\Delta t\,C_m$.
    \item For the instantaneous-quench convention $S_j=1$, including
    $S_0=1$, one obtains
    $F_m^{\mathrm{R}}=\Delta t(C_m-K_m)
    =\Delta t\sum_{k=0}^{m-1}K_k$, exactly the previous exclusive prefix sum.
\end{itemize}
The corrected result satisfies $F_0^{\mathrm{R}}=0$ in both cases.

Alternatively, trapezoidal quadrature assigns half weight to both endpoints:
\begin{equation}
    \boxed{
    F_m^{\mathrm{trap}}=
    \Delta t\left[C_m-\frac{K_0S_m+S_0K_m}{2}\right].
    }
    \label{eq:convolution_trapezoidal}
\end{equation}
This expression also gives $F_0^{\mathrm{trap}}=0$. The endpoint corrections
require only scalar multiplications and additions of MPSs; they do not
require changing the convolution algorithm. The factor $\Delta t$ is applied
once to the output, not separately to each input.

\paragraph{Convolution using the addition tensor.}

Following Section 10.3 of
\textit{Who can compete with quantum computers?},
we introduce the addition tensor
\begin{equation}
    \Theta_{mjk}=\delta_{m,j+k},
\end{equation}
where all indices belong to $\{0,\ldots,L-1\}$.

The discrete convolution can then be expressed as
\begin{equation}
    \boxed{
    F_m \approx \Delta t
    \sum_{j,k=0}^{L-1}
    \Theta_{mjk}S_jK_k.
    }
\end{equation}

The Kronecker delta enforces $m=j+k$, so that
only terms satisfying $j\leq m$ contribute.

The addition tensor admits a low-rank
matrix product operator representation.
By contracting this tensor with the two
quantics MPS representing $S$ and $K$,
we obtain the MPS representation of $F$.

\paragraph{Output grid and boundary conditions.}

Since $S$, $K$, and $F$ each contain $L=2^N$
samples, all three can be represented
using $N$ binary indices.

The physical grids satisfy
\begin{equation}
    \boxed{
    \begin{aligned}
        \mathcal{G}_S &= [T_0,T_M[,\\
        \mathcal{G}_K &= [0,T_M-T_0[,\\
        \mathcal{G}_F &= [T_0,T_M[.
    \end{aligned}
    }
\end{equation}

All three grids have the same spacing $\Delta t$.

The addition MPO must impose a zero
final carry, thereby excluding overflow
and preventing periodic wraparound.

Consequently, no zero-padding is required
to calculate the causal convolution on
the original output grid.

The tensor-network contraction implements
the discrete convolution exactly in the
absence of MPS truncation. The remaining
numerical errors arise from the quadrature
approximation and possible MPS compression.

\paragraph{Special case $T_0=0$.}

When $T_0=0$, the three grids coincide:
\begin{equation}
    \boxed{
    \mathcal{G}_S=\mathcal{G}_K=\mathcal{G}_F
    =[0,T_M[.
    }
\end{equation}

\subsection{Evaluating the integrated response}

The convolution output already represents the integral up to the observation
time:
\begin{equation}
    F(t)=\int_{T_0}^{t}K(t-v)S(v)\,dv.
\end{equation}
Once the convolution, grid-spacing factor, and chosen endpoint treatment
have been applied, evaluating this output MPS at $t$ gives the desired
integral. No second integration is needed. With $F(T_0)=0$, subtracting
$F(T_0)$ is unnecessary; it simply gives $F(t)-F(T_0)=F(t)$.
The evaluation time must belong to the represented grid: the upper boundary
$T_M$ is excluded and requires separate endpoint handling if needed.

This output is not, in general, a primitive for integration over arbitrary
subintervals with a fixed observation time. Such a primitive would be
\begin{equation}
    P_t(x)=\int_{T_0}^{x}K(t-v)S(v)\,dv,
\end{equation}
with $t$ held fixed while $x$ varies. Although $P_t(t)=F(t)$, generally
$P_t(x)\neq F(x)$, since the kernel in $F(x)$ is $K(x-v)$.
Consequently,
\begin{equation}
    F(b)-F(a)\neq\int_a^b K(b-v)S(v)\,dv
\end{equation}
in general. This does not prevent using the convolution output directly
for the intended integral from the initial time to the observation time.

\subsection{Application to the first-order Green function}

For the relative-time convention $v=t-t_1\geq0$, the required contribution is
\begin{equation}
    G^{(1)}(t=t')=
    iU\int_0^{t-T_0}\lambda(t-v)
    \widetilde G^{(1)}(0,0;-v)\,dv.
\end{equation}
To distinguish the roles of the two inputs from the earlier notation,
write the saved integrand as the signal on the positive-delay grid:
\begin{equation}
    S_{\mathrm{GF}}(v)=i\widetilde G^{(1)}(0,0;-v).
\end{equation}
The saved \texttt{t\_ti} integrand MPS already contains both the negative
internal-time argument and the prefactor $i$. Neither an additional
reflection nor another multiplication by $i$ is required.

Introduce elapsed observation time $T=t-T_0$ and the activation input
\begin{equation}
    A(w)=\lambda(T_0+w),\qquad w\geq0.
\end{equation}
Then the desired response is a convolution of these two inputs:
\begin{equation}
    \boxed{
    R(T)=\int_0^T S_{\mathrm{GF}}(v)A(T-v)\,dv,
    \qquad G^{(1)}(t=t')=U R(t-T_0).
    }
\end{equation}
Both input arrays are indexed by elapsed coordinates on
$[0,T_M-T_0[$, with the same number of points and spacing $\Delta t$.
Thus, after applying the quadrature convention, a helper such as
\texttt{make\_indefinite\_integral\_conv} is intended to construct
\begin{equation}
    R_m\approx\Delta t\sum_{j=0}^{m}
    S_{\mathrm{GF}}(j\Delta t)A((m-j)\Delta t).
\end{equation}
For an adiabatic ramp with $A(0)=\lambda(T_0)=0$, this inclusive sum
already equals the exclusive rectangle sum over lag indices $j<m$;
its diagonal endpoint term vanishes and $R_0=0$.

The final calculation therefore consists only of evaluating the response
MPS and multiplying by $U$. If its returned grid represents elapsed time,
evaluate it at $t-T_0$. If the same array is assigned the physical-time
grid $[T_0,T_M[$ instead, evaluate it at $t$. For $T_0=0$, these coordinates
coincide. The grid-spacing factor is applied during response construction,
not again during evaluation, and there is no need to subtract the origin
value for the ramp with $\lambda(T_0)=0$.

\end{document}
